\documentclass[11pt]{article}
\usepackage[T1]{fontenc}
\usepackage{amsmath,amssymb,amsthm}
\usepackage[margin=1in]{geometry}
\usepackage[hidelinks]{hyperref}
\emergencystretch=3em

\newcommand{\D}{\mathbb D}
\newcommand{\Hmat}{\mathcal H}

\title{Literature-implied partial answer for arXiv:2201.09591}
\author{agent\_super\_00}
\date{2026-07-01}

\begin{document}
\maketitle

\paragraph{Status.}
\texttt{literature\_implied\_answer (partial subcase)}.
This note records a later literature theorem, not a new proof by this run.

\paragraph{Source/open problem.}
The source paper is Mikael Lindstrom, Santeri Miihkinen, and David Norrbo,
\emph{Exact essential norm of generalized Hilbert matrix operators on classical
analytic function spaces}, arXiv:2201.09591v1 (2022).
In Section 9, ``Behavior of norm and essential norm of the Hilbert matrix
operator'', page 22 of the arXiv PDF, the authors state that it is still open
to determine the exact norm of the Hilbert matrix operator $\Hmat$ on standard
weighted Bergman spaces in the range
\[
  2+\alpha < p <
  2+\alpha+\sqrt{(2+\alpha)^2-\left(\sqrt 2-\frac12\right)(2+\alpha)}.
\]
The surrounding paragraph explains why equality of norm and essential norm is
relevant; Theorem 9.3 and Corollary 9.4 show that nonattainment of the norm
would imply
\[
  \|\Hmat\|_{A^p_\alpha\to A^p_\alpha}
  = \|\Hmat\|_{e,A^p_\alpha\to A^p_\alpha}
  = \frac{\pi}{\sin((2+\alpha)\pi/p)}.
\]

\paragraph{Supporting theorem.}
The supporting paper is Guanlong Bao, Liu Tian, and Hasi Wulan,
\emph{The norm of the Hilbert matrix operator on Bergman spaces},
arXiv:2601.13672v1 (2026).
Its abstract says that it proves Karapetrovic's conjecture for the range
\[
  0\le \alpha \le
  \frac{6p^3-29p^2+17p-2+2p\sqrt{6p^2-11p+4}}{(3p-1)^2}.
\]
Theorem 4.1, labeled \texttt{corr} in the TeX source and appearing on page 15
of the arXiv PDF, states precisely that under this hypothesis
\[
  \|\Hmat\|_{A^p_\alpha}
  = \frac{\pi}{\sin((\alpha+2)\pi/p)}.
\]

\paragraph{Identification and scope.}
The source problem is the exact norm problem for the same Hilbert matrix
operator on the same standard weighted Bergman spaces. Thus Theorem 4.1 of
arXiv:2601.13672 answers the source open problem on the intersection of the
source's open interval with the displayed 2026 parameter range. This is only a
partial answer. The supporting paper itself states in the introduction that the
conjecture was not fully settled for $\alpha>0$ and
$\alpha+2<p<2(\alpha+2)$, and that the negative-indexed region
$-1<\alpha<0$ remains open in the corresponding strip.

\paragraph{Provenance.}
The supporting paper cites the source paper as Lindstrom--Miihkinen--Norrbo,
but it does not explicitly say that Theorem 4.1 is an answer to the page-22
open-problem sentence of arXiv:2201.09591. The connection is therefore an
agent-identified literature implication, not an explicitly announced later
answer.

\paragraph{Files and checks.}
The packet includes local copies of arXiv:2201.09591 and arXiv:2601.13672.
The identification was checked against the cached TeX sources for
arXiv:2201.09591 and arXiv:2601.13672, plus the run's cheap indexes.

\begin{thebibliography}{9}
\bibitem{LMN}
M. Lindstrom, S. Miihkinen, and D. Norrbo,
\emph{Exact essential norm of generalized Hilbert matrix operators on
classical analytic function spaces}, arXiv:2201.09591v1, 2022.

\bibitem{BTW}
G. Bao, L. Tian, and H. Wulan,
\emph{The norm of the Hilbert matrix operator on Bergman spaces},
arXiv:2601.13672v1, 2026.
\end{thebibliography}

\end{document}
